Menciona los tipos de protocolos de red.

Los protocolos de red son conjuntos de reglas y estándares formales que permiten a dos o más dispositivos comunicarse e intercambiar información a través de una red. Se pueden clasificar como sigue~\cite{tanenbaum_redes}:
\begin{itemizeColor}
    \item \textbf{Protocolos de comunicación:} Se encargan del direccionamiento lógico y del transporte de los datos entre dispositivos finales. Incluyen protocolos no orientados a conexión como IP (IPv4 e IPv6) y UDP, y protocolos orientados a conexión y confiables como TCP.
    \item \textbf{Protocolos de enrutamiento:} Descubren la topología de la red y calculan las rutas óptimas para reenviar los paquetes entre distintos sistemas autónomos o subredes (por ejemplo, BGP, OSPF y RIP).
    \item \textbf{Protocolos de gestión, control y soporte:} Permiten diagnosticar la conectividad, resolver direcciones físicas o asignar configuraciones dinámicas de red de forma automática. 
    \item \textbf{Protocolos de seguridad:} Proveen mecanismos de cifrado, autenticación e integridad para proteger los datos en tránsito frente a intercepciones y modificaciones.
    \item \textbf{Protocolos de aplicación:} Definen los formatos y reglas de intercambio para servicios específicos dirigidos al usuario final o a los procesos del sistema.
\end{itemizeColor}

